% demo.tex — minimal, self-contained example of the Triton LaTeX integration.
%
% Build the figures first, then compile this file:
%
%   make figures        # ../dist/cli.cjs render-dir diagrams/ -o figures/
%   make pdf            # tectonic demo.tex   (or pdflatex demo.tex)
%
% On Overleaf: just upload this folder. The figures/*.pdf are committed, so no
% Node/Triton toolchain is needed — graphicx includes the vector PDFs directly.

\documentclass{article}
\usepackage[margin=1in]{geometry}
\usepackage{triton}

% Point at the committed, precompiled vector PDFs.
\tritondir{figures}

\title{Triton --- LaTeX Integration Demo}
\author{The \texttt{@triton/latex} package}
\date{}

\begin{document}
\maketitle

Triton diagrams are precompiled to \emph{vector} PDF and dropped in by name with
\verb|\triton{<name>}|. The assets below are committed to the repository, so this
document compiles on pdf\LaTeX{}, Xe\LaTeX{}, Lua\LaTeX{}, and Overleaf with no
external toolchain.

\section{Inline include}

A flowchart, scaled to the text width:

\begin{center}
  \triton{flowchart}
\end{center}

\section{Figure with caption}

The \verb|\tritonfig| helper wraps a diagram in a \texttt{figure} with a caption:

\tritonfig[width=0.5\linewidth]{avl}{An AVL tree rendered by Triton.}

\section{A composed poster}

Span and scale like any other graphic:

\begin{center}
  \triton[width=\linewidth]{ds-poster}
\end{center}

\end{document}
